\documentclass[10pt,a4paper]{amsart}
\usepackage[margin=19mm]{geometry}
\usepackage{amsmath,amssymb,mathtools}
\usepackage[scr=rsfs,cal=euler]{mathalfa}
\usepackage{xfrac}
\usepackage[unicode,hidelinks,bookmarks=false]{hyperref}
\newcommand{\ie}{i.e.\ }
\newcommand{\cf}{cf.\ }
\newcommand{\resp}{respectively\ }
\newcommand{\Subsection}{\subsection}
\providecommand{\nobreakdash}{\nobreak\hbox{-}}
\setlength{\emergencystretch}{2em}
\multlinegap0pt
\usepackage{bm}
\usepackage{bbm}
\usepackage{dsfont}
\usepackage{xspace}
\usepackage{cancel}
\usepackage{mathtools}
\usepackage{amsthm}
\makeatletter
\@ifundefined{theorem}{\newtheorem{theorem}{Theorem}}{}
\@ifundefined{lemma}{\newtheorem{lemma}[theorem]{Lemma}}{}
\makeatother
\newcommand{\Prob}{\mathbb{P}}
\title[Dense Weak Hiding: Closing Complexity Gaps in Nonconvex and ]{Dense Weak Hiding: Closing Complexity Gaps in Nonconvex and PL Finite-Sum Optimization under Individual Smoothness}
\author{Yuxing Peng, Zhiqing Tang, Weijia Jia}
\date{Source: arXiv:2609.00045v2 (CC BY 4.0)}
\begin{document}
\maketitle

\noindent\emph{Source and licence.} Dense Weak Hiding: Closing Complexity Gaps in Nonconvex and PL Finite-Sum Optimization under Individual Smoothness. Version arXiv:2609.00045v2 (CC BY 4.0), distributed under the Creative Commons Attribution 4.0 International licence, as recorded in the arXiv licence metadata for this exact version and shown on the abstract page https://arxiv.org/abs/2609.00045v2. Full rights notice: licenses/arxiv-2609.00045-CC-BY-4.0.txt.\par\smallskip

\noindent\textit{Editorial scope.} Counted unit: the Lemma whose source label is \texttt{lem:first-use-row-information} in the article (source0.tex lines 1654--1673, source label \texttt{lem:first-use-row-information}), delivered together with the article's own complete original proof of it (source0.tex lines 1675--1812, one \texttt{proof} environment of 138 lines). No mathematics is written, repaired or re-derived: statement and proof are the article's own text line for line. Required context retained uncounted: none. Every definition, display and label of those lines is retained unchanged. Omitted from this package and not redistributed: 1--1653, 1674--1674, 1813--5349. Nothing inside a retained excerpt is omitted or altered: every retained body is delivered exactly as the article's lines are written, so no omission receipt of any kind accompanies this package. Citations: the retained excerpts carry no citation command at all, so no bibliography entry of the article is transcribed and no reference is redistributed. Reference adaptation: none was needed and none was made; every reference inside the delivered unit resolves inside it, so no unresolved reference or citation is produced. Printed slips kept literal: no printed word, symbol, hypothesis or number of the counted statement or of its proof was altered. Layout and class: the article's own class is \texttt{article} with packages including iclr2027\_conference, times, amsmath, amssymb, amsthm, mathtools, booktabs, xcolor, colortbl, hhline, tabularx, multirow, float, algorithm; the wrapper is the lane's pinned \texttt{amsart} editorial wrapper at 10pt on A4, which restates the theorem-like environments the retained text needs and adds the article's own macro definitions that the retained text needs (\texttt{\textbackslash Prob}), with extra wrapper packages (bm, bbm, dsfont, xspace, cancel, mathtools). The delivered numbering is the unit's own. Assets: no retained span contains a figure environment, a graphics-inclusion command, a PDF-page inclusion, a table or a TikZ picture; no image file is copied into this package, the package declares no asset dependency and no image file is redistributed with it. Licence: version arXiv:2609.00045v2 of this article is distributed under the Creative Commons Attribution 4.0 International licence, as recorded in the arXiv licence metadata for this exact version and shown on the abstract page https://arxiv.org/abs/2609.00045v2. A full rights notice is delivered at licenses/arxiv-2609.00045-CC-BY-4.0.txt. Characters: every retained excerpt is pure ASCII, so the delivered engine, XeLaTeX with its default Latin Modern text font, reports no missing character.
\begin{lemma}[Information in the first $r$ distinct rows]
\label{lem:first-use-row-information}
Let $1\le r<n$ and assume
\begin{equation}
 8m+9r\le n.
 \label{eq:row-information-interior}
\end{equation}
For the stopped history $\mathcal T_r$,
\begin{equation}
 \mathrm I(\Theta;\mathcal T_r)
 \le
 \frac{32D}{63}
 \sum_{d=0}^{r-1}
 \left(\frac{m}{n-d}\right)^2
 \le
 \frac{32D}{63}\,
 r\left(\frac{m}{n-r}\right)^2.
 \label{eq:first-use-row-information}
\end{equation}
\end{lemma}

\begin{proof}[Proof of Lemma~\ref{lem:first-use-row-information}]
\emph{Information in a newly queried row.}
Consider a query $t$ retained in $\mathcal T_r$ and condition on a realization
$h$ of $\mathcal H_{t-1}$.  Both $I_t$ and $K_t$ are then fixed.  If $I_t$ is
repeated, its row and the corresponding reply are already determined by $h$,
so the conditional information increment is zero.  Suppose instead that
$I_t$ is unseen and $K_t=d<r$.  For column $\ell$, let $\Sigma_{d,\ell}$ be the
sum of that coordinate over the $d$ distinct rows already revealed in $h$.
Conditional exchangeability leaves $n-d$ uniformly assigned signs with
residual sum $m\Theta_\ell-\Sigma_{d,\ell}$.  Hence
\begin{equation}
 p_{\theta,\ell}^{(d)}
 :=
 \Prob(S_{I_t\ell}=1
 \mid \Theta_\ell=\theta,\mathcal H_{t-1}=h)
 =
 \frac{n+m\theta-d-\Sigma_{d,\ell}}{2(n-d)}.
 \label{eq:conditional-unseen-sign}
\end{equation}
The optional initialization row obeys the same formula with $d=0$, because its
index is fixed by $W$ before the row is observed.

Expose the coordinates of the new row in the order $\ell=1,\ldots,D$.
Conditional independence of the column sequences gives
\[
 \Prob(S_{I_t\ell}=1
 \mid\Theta,\mathcal H_{t-1}=h,S_{I_t,<\ell})
 =p_{\Theta_\ell,\ell}^{(d)}.
\]
Thus, under this conditioning, the current sign depends on the hidden
direction only through $\Theta_\ell$.  The argument does not require
independence among the posterior coordinates of $\Theta$.
More explicitly, conditional on $\mathcal H_{t-1}=h$ and
$S_{I_t,<\ell}$,
\[
 \Theta\longrightarrow\Theta_\ell\longrightarrow S_{I_t\ell}
\]
is a Markov chain.  Therefore,
\[
 \mathrm I(\Theta;S_{I_t\ell}
 \mid\mathcal H_{t-1}=h,S_{I_t,<\ell})
 =
 \mathrm I(\Theta_\ell;S_{I_t\ell}
 \mid\mathcal H_{t-1}=h,S_{I_t,<\ell}).
\]

Condition~\eqref{eq:row-information-interior} implies
\[
 \left|p_{\theta,\ell}^{(d)}-\frac12\right|
 \le
 \frac{m+d}{2(n-d)}
 \le\frac1{16}.
\]
Hence both $p_{+,\ell}^{(d)}$ and $p_{-,\ell}^{(d)}$ lie in
$[7/16,9/16]$.  For $Y\sim\operatorname{Bernoulli}(p)$, its Shannon entropy is
$H(Y)=h_{\rm b}(p):=-p\log p-(1-p)\log(1-p)$.  Here $\log$ is the natural
logarithm, so entropy and mutual information are measured in nats.  On the
preceding interval,
\[
 -h_{\rm b}''(p)=\frac1{p(1-p)}
 \le\frac{256}{63}
 \qquad (p\in[7/16,9/16]).
\]

Fix also a realization of $S_{I_t,<\ell}$ and write
\[
 \pi:=\Prob(\Theta_\ell=1\mid
   \mathcal H_{t-1}=h,S_{I_t,<\ell}),
 \qquad
 \bar p:=\pi p_{+,\ell}^{(d)}+(1-\pi)p_{-,\ell}^{(d)}.
\]
Under this conditioning, $S_{I_t\ell}$ is Bernoulli with parameter $\bar p$;
after additionally conditioning on $\Theta_\ell=\theta$, its parameter is
$p_{\theta,\ell}^{(d)}$.  Thus the first conditional entropy below is
$h_{\rm b}(\bar p)$, while the entropy after conditioning on $\Theta_\ell$ is
the posterior average of $h_{\rm b}(p_{+,\ell}^{(d)})$ and
$h_{\rm b}(p_{-,\ell}^{(d)})$.  The identity
$\mathrm I(X;Y\mid Z)=H(Y\mid Z)-H(Y\mid X,Z)$ gives
\begin{align}
 \mathrm I(
 \Theta;S_{I_t\ell}
 \mid \mathcal H_{t-1}=h,S_{I_t,<\ell})
 &=h_{\rm b}(\bar p)
   -\pi h_{\rm b}\!\left(p_{+,\ell}^{(d)}\right)
   -(1-\pi)h_{\rm b}\!\left(p_{-,\ell}^{(d)}\right)
   \notag\\
 &\le \frac12\frac{256}{63}\,
   \pi(1-\pi)
   \left(p_{+,\ell}^{(d)}-p_{-,\ell}^{(d)}\right)^2
   \notag\\
 &\le\frac{32}{63}
   \left(\frac{m}{n-d}\right)^2.
 \label{eq:one-coordinate-mutual-information}
\end{align}
The first inequality is Taylor's theorem for the concave function $h_{\rm b}$:
its Jensen gap is at most one half of the curvature bound times the variance of
the two-point parameter.  The last inequality uses
\[
 \pi(1-\pi)\le\frac14,
 \qquad
 p_{+,\ell}^{(d)}-p_{-,\ell}^{(d)}=\frac{m}{n-d}.
\]

\emph{Chain rule along the actual history.}
Because the coordinate bound is uniform in $S_{I_t,<\ell}$, the chain rule
within the new row gives
\begin{align}
 \mathrm I(\Theta;S_{I_t:}\mid\mathcal H_{t-1}=h)
 &=\sum_{\ell=1}^D
   \mathrm I(\Theta;S_{I_t\ell}
   \mid\mathcal H_{t-1}=h,S_{I_t,<\ell})
   \notag\\
 &\le \frac{32D}{63}\left(\frac{m}{n-d}\right)^2.
 \label{eq:one-new-row-information}
\end{align}
Pad the stopped history after its stopping time with deterministic null
query--response symbols to a total of $B$ rounds.  The chain rule then has one
increment in conditional mutual information per round.  The chosen query is
$\mathcal H_{t-1}$-measurable and therefore adds no information by itself; a
response from a repeated row, or a null response, is already determined by the
preceding history and also adds zero.
The optional initialization row, when present, is the $d=0$ increment before
the first query.  Each response from a new row increases the number of distinct rows, so
every $d=0,\ldots,r-1$ contributes at most once.  (Without an initialization
row, the first response from a new row corresponds to the $d=0$ increment.)  Applying
\eqref{eq:one-new-row-information} at these increments gives
\begin{align*}
 \mathrm I(\Theta;\mathcal T_r)
 &\le \frac{32D}{63}
 \sum_{d=0}^{r-1}\left(\frac{m}{n-d}\right)^2
 \le \frac{32D}{63}\,
 r\left(\frac{m}{n-r}\right)^2.
\end{align*}
Conditioning on $S_{I_t,<\ell}$ in
\eqref{eq:one-coordinate-mutual-information} makes posterior independence
unnecessary.  Finally, $n-d\ge n-r$ for $d\le r-1$.  This
proves~\eqref{eq:first-use-row-information}.
\end{proof}

\end{document}
